\section{Threats to validity}\label{sec:threats}

The limitations below are stated twice over: as measurement caveats, and as risks for anyone
relying on the catalogue for the purpose Section~\ref{sec:legal} proposes.

\paragraph{Duration bias, and the bound it places on every claim.}
The corpus admits a case only if it both opened and concluded within the same calendar year.
The admissible lifetime therefore shrinks month by month: a January case has 365 days to
conclude, a December case has 31, and the observed maxima track that ceiling almost exactly
(343 days against 365 in January; 29 against 31 in December). Two losses stack, and they are of
different kinds. Of the 12{,}621 cases opened in the year, 7{,}176 had not concluded by
31 December and are excluded outright; the analysed corpus is 41.5\% of the year's openings.
Within the concluded population, taking the least-truncated month as a reference, a further
18.2\% is estimated missing. Both losses fall on the long cases. Three consequences bind the
paper. The corpus is not
``cases opened in the year'' but ``cases opened in the year that also finished quickly enough'';
no duration or efficiency statistic from it is admissible, and the observed median of 14 days is
biased downward by construction; and the catalogue under-represents long, multi-unit
procedures --- procurement, contracting, payment chains --- which are precisely the types where
routing analysis would have most operational value. The one-year scope was a cost decision for
an initial study rather than a methodological preference. For legal use the bias runs in the
worst available direction: the cases the corpus loses are the slow ones, which are the cases
that generate complaints about duration and the cases a reasonable-duration guarantee exists to
address. A service charter populated from this corpus would understate the expected time, and
would understate it most where the understatement misleads.

\paragraph{Temporal dependency is a scope boundary, not a defect.}
The partition uses each case's complete history by design (Section~\ref{sec:res-ablation}), and
a real share of its signal arrives late. Nothing reported here transfers to assigning a type at
intake, which would require early-stage features the study does not have. An entry says what a
concluded case turned out to be.

\paragraph{All models under-generate.}
Generated paths match roughly 45\% of a real case's visits, and are shorter than the medoid
(11.6--12.6 visits against 13.9). The typical paths capture what recurs, not the full extent of
what a case does. Published as a statement of the main processing steps, a path that omits steps
understates what a user must go through, and the omission is systematic rather than random.

\paragraph{The internal-validity result is underpowered.}
Sixteen clusterings; no index reaches $p < 0.05$ under permutation. The finding is absence of
evidence, and a properly powered replication is future work.

\paragraph{Validity indices are computed on noise-free subsets.}
The internal indices mask unassigned cases, so density-based clustering is scored on the
83--86\% it retained while the other algorithms are scored on everything. Neither the intrinsic
nor the downstream comparison is valid without correcting for this, which
Section~\ref{sec:res-rq2} does for the latter.

\paragraph{Cluster stability is a lower bound.}
The clusterer is refit inside the bootstrap; the dimensionality reduction is not, so the
reported Jaccard recovery understates total variability. It is also a subsample estimate over
25 draws, and it exists only for the two density-based clusterings.

\paragraph{Indices computed on a reduced embedding.}
Clustering and validity are both computed in the reduced space. This supports comparison between
methods on the same embedding rather than absolute claims about structure in the original
feature space. The concern is documented for two dimensions and we extend it cautiously to
fifty.

\paragraph{Only the opening document is embedded.}
The document pipeline ran over every case opened in the year, embedding 12{,}259 opening
documents --- one per case at that scope --- of which the 5{,}243 belonging to analysed cases
enter the feature matrix. Against 148{,}178 documents held for those cases, the semantic block
is therefore temporally safe by construction but sees a small fraction of the documentary
record.

\paragraph{Small clusters dominate the counts.}
278 of the evaluated clusters on the action alphabet fall in the smallest band, evaluated on
about eleven held-out cases each. All size-dependent results are reported stratified for this
reason, and the pooled averages should not be read alone. Entries differ in how well they are
evidenced, which is why case counts travel with each type in the catalogue.

\paragraph{A strong baseline we do not dominate, and beat on one alphabet only.}
The medoid outperforms the generated paths on several support measures, and on the action
alphabet it also fabricates less per visit than every generator --- by a factor of 1.8 against
the best of them. The fabrication argument for generated paths therefore holds for the routing
representation and not for the action one (Section~\ref{sec:res-rq3}). The medoid is also
computed against a sample of at most twenty training traces rather than the whole cluster, so
on large clusters the baseline reported here is weaker than a full-cluster medoid would be, and
the comparison understates it.

\paragraph{The motivating premise is imported, not tested here.}
That transparency regimes fail because the record was never created is established for a
different jurisdiction, disclosure regime and subject matter \citep{brauneis2018algorithmic}.
We did not establish it for the administration studied (Section~\ref{sec:legal}). The method
stands or falls on its own evidence either way, but the argument that the artefact is
\emph{needed} rests on an extrapolation this study does not close.

\paragraph{Single administration, single system.}
One state secretariat, one electronic process system, one year. Whether the catalogue structure
or the size threshold transfers to other administrations is untested. The normative claim is
bounded no less tightly: it runs where a proactive duty to publish processing steps in
machine-readable form exists, as Section~\ref{sec:legal} sets out for this jurisdiction. Where
disclosure is request-driven instead, the case for producing the artefact unprompted must be
made on other grounds.

\paragraph{The types and their names have not been assessed by specialists.}
Names and typical paths were produced by the pipeline and read by the authors; the in-venue
precedent validates such structure with domain specialists \citep{vianna2024topic}, and no
comparable study is reported here. The names deserve separate mention from the types. They are
generated by a language model from a sample of each cluster's cases (Section~\ref{sec:naming}),
in Portuguese, over administrative subject matter, and no domain expert has checked that any of
them describes what the cluster contains. Since the name is what a reader of a published
catalogue encounters first and often only, an unvalidated name is the component most likely to
mislead, and near-duplicate naming was an observed failure mode during development. Until an
assessment exists, an entry is a hypothesis evidenced by the cases assigned to it, not a
description on which a decision about a case may rest.

\paragraph{Parallel routing, linearised.}
A case can run in more than one organisational unit at once, and the routing alphabet represents
it as a single sequence ordered by event timestamp (Section~\ref{sec:alphabets}). Concurrent
stretches therefore appear as though they were successive. Two consequences follow. An observed
transition between two units is not always a handover from one to the other: it may be only the
order in which two parallel stretches happened to record their events. And the typical paths
inherit the simplification --- they describe a plausible order of visits rather than the
procedure's concurrency structure --- so nothing here about precedence between units should be
read as a claim of dependency between them. Modelling concurrency explicitly would require a
trace representation other than a sequence, and is future work.

\paragraph{Uncollectable documents.}
96 opening documents were cancelled in the source system and can never be retrieved; 202 cases
have no embedded document and are excluded from clustering.
